\documentclass{subfiles}
\begin{document}
    Above we've used the term information assuming its meaning as given; 
        it should be clear that the notion of information in the spoken language
        has little to do with the CS counterpart. 

    Formally, the notion of information is that of a measure of the
        reduction of the uncertainty of the outcomes of an event. 
        That is, a measure of how much knowledge is gained observing the event. 
        Therefore, it makes sense that the information carried by an event is
        inversely proportional to its probability.

    Historically, the first person attempting to define information was R.V.L. Hartley;
        see \cite{hartley1928} for more details. But the major contribution is due to 
        Claude Shannon, who in \cite{shannon1948}, building on Hartley's ideas,
        provides a rigorous definition of information and a way to compute it.

    In his paper, Shannon shows that any communication system can be reduced to 
        five parts: a \emph{sender}, a \emph{transmitter}, a \emph{communication channel}, a \emph{receiver} and a \emph{destination}. 
        A schematic representation of such system is shown in \Cref{fig:shannonSystem}.

    With his paper, Shannon introduces introduces two important hypotesis:
        \begin{itemize}
            \item Transmitting a symbol requires a non-zero amount of time;
            \item The communication system is subject to noise.
        \end{itemize}

        \subfile{../../TikZ/Figure 1 - Shannon communication system}

    \begin{remark*}
        In his paper Shannon actually develops two theories: 
            one with noise and one without. We focus on the latter.
    \end{remark*}

    The question now is: what properties should a measure of information satisfy?
        According to Shannon, such a function should have a dependence on probability,
        should be monotonic, continous and additive for indenpendent events.

    More often than not, we are interestend in computing the information a source;
        that is, an entity over some alphabet \(\Sigma\),
        that at any given time produces a symbol \(\sigma \in \Sigma\)
        with probability \(\Prob{\sigma}\). 
        For instance, a Markov chain is an example of source of information.
        We shall say that there exist two types of sources: with memory and memoryless.
        In what follows, we consider memoryless sources.

    But how do we actually measure information? According to Shannon, 
        the only valid measure is 
        \begin{equation}\label{eq:shannonEntropy}
            H(S) = K \sum_{i = 1}^{n} {p_{i} \log_{b} p_{i}}
        \end{equation}
        where \(n\) is the total number of symbols produced by \(S, p_{i}\)
        the probability associated with the \(i\)-th symbol produced, and \(K\) a 
        positive constant usually set to 1.
        Unless stated otherwise, the base \(b\) of the logarithm is always 2.
        We usually refer to \Cref{eq:shannonEntropy} as \emph{entropy of a source} \(S\).
    
    Although we could prove that \Cref{eq:shannonEntropy} is the only valid measure,
        it is beyond our scope; the interested reader is referred to \cite{shannon1948}[Appendix 2].
\end{document}
